\section{Conclusion}

Aggregate TFP did not grow in Colombia between 2005 and 2024. The chained index fell from 100 to 98.5, a cumulative decline of 1.5 percent and an annualized log change of $-0.081$ percentage points. Most of the decline occurred before 2011, but the subsequent improvement did not become sustained aggregate TFP growth.

The aggregate result was produced by large contributions of opposite signs. Social services, trade, agriculture, and transport contributed 0.361 percentage points per year. Finance and real estate, mining, construction, manufacturing, and electricity subtracted 0.441 points. Finance and real estate and mining made the largest negative contributions, while social services made the largest positive contribution.

The year-by-year aggregation is central to this conclusion. Cumulative sectoral TFP growth does not add across activities, and average weights applied to long-run sectoral growth do not reproduce the exact total when weights change over time. Recovering the annual weights and averaging annual contributions preserves the published aggregate identity.

The accounting counterfactual shows the scale of the negative contributions. If annual TFP growth in the four weakest activities had been zero, production-account output growth would have been 3.59 percent per year rather than 3.15 percent. The 0.43-point difference is economically meaningful, but it is not a causal estimate or a forecast of GDP. The result identifies where a more detailed diagnosis may have the highest aggregate payoff. It does not identify the policies that would produce those gains.
